\section{Производные для регулярных языков}

Производные для регулярных языков предложены в Янушем Бжозовским в работе <<Derivatives of Regular Expressions>> ~\sidecite{Brzozowski1964}.
Будучи достаточно удобным механизмом для решения ряда (в том числе прикладных) задач, производные активно исследовались и мы будем использовать работу~\sidecite{OWENS_REPPY_TURON_2009} как основу для нашего изложения.

Общее определение производных было дано выше и работает, разумеется, и для регулярных языков.
Однако общее определение ничего не говорит о том, как конструктивно вычислить производную.
Для регулярных языков возможно представить алгоритм, который вычисляет производные на основе регулярных выражений.

Прежде чем определить непосредственно производные, определим вспомогательную функцию, которая будет по регулярному выражению определять, принадлежит ли пустая цепочка языку, задаваемому им.
Часто эта функция называется \emph{IsNull}.

\begin{itemize}
    \item $\emph{IsNull}(\varepsilon) = \emph{true}$
    \item $\emph{IsNull}(\varnothing) = \emph{false}$
    \item $\emph{IsNull}(x) = \emph{false}$
    \item $\emph{IsNull}(R_1 \cdot R_2) = \emph{IsNull}(R_1) \ \& \ \emph{IsNull}(R_2) $
    \item $\emph{IsNull}(R_1 \mid R_2) = \emph{IsNull}(R_1) \mid \emph{IsNull}(R_2) $
    \item $\emph{IsNull}(R^*) = \emph{true}$
\end{itemize}

Далее, введём функцию $\mathcal{V}$:
\[
    \mathcal{V}(R) =
    \begin{cases}
        \varepsilon, \emph{IsNull}(R) \\
        \varnothing, \text{ иначе}
    \end{cases}.
\]

Теперь мы готовы определить производные для различных конструкций регулярных выражений.
\begin{itemize}
    \item $\partial_t(\varepsilon) = \varnothing$
    \item $\partial_t(\varnothing) = \varnothing$
    \item $\partial_t(t) = \varepsilon$
    \item $\partial_t(x) = \varnothing \text{ для } t \neq x$
    \item $\partial_t(R_1 \cdot R_2) = \partial_t(R_1) \cdot (R_2) \mid \mathcal{V}(R_1) \cdot \partial_t(R_2)$
    \item $\partial_t(R_1 \mid R_2) = \partial_t(R_1) \mid \partial_t(R_2)$
    \item $\partial_t(R^*) =$\sidenote{Интересное упражнение~--- показать это, расписав по определению звезду Клини.} $ \partial_t(R)\cdot R^*$.
\end{itemize}

\begin{example}
    Вычислим производную $a^*(a \mid b)$ по символу $a$.
    \[
        \partial_a(a^*(a \mid b)) = \partial_a(a^*) (a \mid b) \mid \mathcal{V}(a^*) \cdot \partial_a(a \mid b) \\
    \]
    Вычислим отдельные подвыражения.
    \begin{align*}
        \partial_a(a^*) & = \partial_a(a) a^* = a^* \\
        \partial_a(a \mid b) & = \partial_a(a) \mid \partial_a(b) = \varepsilon \mid \varnothing = \varepsilon
    \end{align*}
    Соединим всё вместе.
    \begin{align*}
        \partial_a(a^*(a \mid b)) & = \partial_a(a^*) (a \mid b) \mid \mathcal{V}(a^*) \cdot \partial_a(a \mid b) \\
                                  & = a^* (a \mid b) \mid \varepsilon \cdot \varepsilon \\
                                  & = a^* (a \mid b) \mid \varepsilon \\
                                  & = (a^* (a \mid b))?
    \end{align*}
\end{example}

Одной из прикладных задач, которую можно решить с помощью производных, является принадлежность строки языку, задаваемому регулярным выражением\sidenote{Что важно, без преобразований регулярного выражения в конечный автомат, как было показано выше.}.
Пусть дано регулярное выражение $R$ и нужно проверить, принадлежит ли строка $w=t_0 \ldots \ t_{n_1}$ языку, задаваемому этим выражением.
Для этого необходимо вычислить следующее выражение:
\[
\emph{IsNull}(\partial_{t_{n-1}}(\ldots(\partial_{t_0}(R))\ldots)).
\]
Если его значение \emph{true}, то строка принадлежит языку.
Иначе не принадлежит.

\begin{example}
    Пусть дано регулярное выражение $a(a \mid (b \ b))^*$ и строка \texttt{abb}.
    Принадлежит ли она языку, задаваемому данным выражением?
    Для ответа на этот вопрос вычислим следующее выражение:
    \[
        \emph{IsNull}(\partial_b(\partial_b(\partial_a(a^*(a \mid (b \ b))^*)))).
    \]
    Для начала вычислим производные.
    \begin{align*}
        \partial_a(a^*(a \mid (b \ b))^*) &= \partial_a(a) (a \mid (b \ b))^* \mid \mathcal{V}(a)(\partial_a(a \mid (b \ b))^*) \\
                                          &= \varepsilon (a \mid (b \ b))^* \mid \varnothing (\partial_a(a \mid (b \ b))^*) \\
                                          &= (a \mid (b \ b))^* \mid \varnothing \\
                                          &= (a \mid (b \ b))^*\\
        \partial_b((a \mid (b \ b))^*)    &= \partial_b((a \mid (b \ b))) (a \mid (b \ b))^* \\
                                          &= (\varnothing \mid b) (a \mid (b \ b))^* \\
                                          &= b (a \mid (b \ b))^* \\
        \partial_b(b (a \mid (b \ b))^*)  &= (a \mid (b \ b))^*
    \end{align*}
    $\emph{IsNull}((a \mid (b \ b))^*) = \emph{true}$ по определению.
    Таким образом, цепочка \texttt{abb} принадлежит языку, задаваемому регулярным выражением $a(a \mid (b \ b))^*$.
\end{example}
